\documentclass[12pt,a4paper]{article}
\usepackage{iftex}
\ifPDFTeX
    \usepackage[utf8]{inputenc}
    \usepackage[T1]{fontenc}
    \usepackage{mathpazo} % Palatino font
\else
    \usepackage{fontspec}
    \setmainfont{Times New Roman}
\fi
\usepackage[brazilian]{babel}
\usepackage{amsmath, amssymb, amsthm}
\usepackage[linesnumbered,ruled,vlined,portuguese]{algorithm2e}
\usepackage{listings}
\usepackage{booktabs}
\usepackage{geometry}
\usepackage{float}
\usepackage{xcolor}
\usepackage{setspace}
\usepackage[colorlinks=true, linkcolor=blue, urlcolor=blue, citecolor=blue]{hyperref}

\geometry{margin=2.5cm}
\onehalfspacing

% Colors for code blocks
\definecolor{codebg}{rgb}{0.96,0.96,0.96}
\definecolor{codekey}{rgb}{0.0,0.4,0.8}
\definecolor{codestr}{rgb}{0.6,0.1,0.1}
\definecolor{codecom}{rgb}{0.4,0.6,0.4}

\lstset{
    backgroundcolor=\color{codebg},
    commentstyle=\itshape\color{codecom},
    keywordstyle=\bfseries\color{codekey},
    stringstyle=\color{codestr},
    basicstyle=\ttfamily\footnotesize,
    breaklines=true,
    numbers=left,
    numberstyle=\tiny\color{gray},
    frame=single,
    rulecolor=\color{lightgray},
    captionpos=b,
    tabsize=4
}

\title{
    \vspace{-1cm}
    \begin{center}
        \Large \textbf{Universidade Federal do Ceará -- UFC} \\
        \large Departamento de Computação \\
        \vspace{1cm}
        \Huge \textbf{Métodos Numéricos II} \\
        \vspace{0.5cm}
        \Large \textit{Relatório da Unidade 2: Integração Numérica}
    \end{center}
}
\author{
    Nicolas Wallace Amorim Perote -- Matrícula: 566512 \\
    Samyra Vitória Lima de Almeida -- Matrícula: 521240
}
\date{\today}

\begin{document}

\maketitle
\tableofcontents
\newpage

\section{Introdução}
O cálculo de integrais definidas do tipo $I = \int_a^b f(x)\,dx$ constitui uma das tarefas mais onipresentes na engenharia, física computacional e ciência de dados. Contudo, em inúmeros contextos, a primitiva da função integranda é intratável, ou a função possui pontos singulares de complexa avaliação, exigindo abordagens algorítmicas de aproximação iterativa.

O presente documento expõe as atividades da \textbf{Unidade 2} e engloba o desenvolvimento de uma suíte completa e escalável de integração numérica em \textit{Python}, utilizando vetorialização nativa com \texttt{NumPy}. Foram contemplados os tradicionais métodos baseados em pontos igualmente espaçados (Newton-Cotes), malhas adaptativas geradas por raízes de polinômios ortogonais (Quadraturas de Gauss), transformações hiperbólicas para controle de polos (Integração Exponencial) e a generalização do Teorema de Fubini em malhas $N$-dimensionais (Integrais Duplas).

A implementação foi estruturada sob os preceitos de Orientação a Objetos, garantindo encapsulamento, generalização (por exemplo, na aplicação automática das regras \textit{compostas}) e extrema clareza semântica, uma evolução significativa se comparada a abordagens em múltiplos \textit{scripts} procedurais.

\section{Fórmulas de Newton-Cotes}
As fórmulas de Newton-Cotes resolvem a integração de $f(x)$ substituindo a função por um polinômio de grau $n$ (via interpolação de Lagrange) em nós perfeitamente eqüidistantes no domínio $[a, b]$. A suíte implementa nativamente o suporte para \textbf{intervalos compostos}: o domínio inicial é subdividido e, em cada subseção, as regras simples (Trapézio, Simpson, Boole, etc.) são aplicadas localmente.

Diferenciamos as metodologias entre \textit{fechadas} (onde os limites superior e inferior também são nós de avaliação) e \textit{abertas} (avaliando apenas os pontos interiores, algo especialmente profícuo quando existem indeterminações algébricas em $x=a$ ou $x=b$).

\begin{lstlisting}[language=Python, caption={Extrato: Suporte a fórmulas compostas de Newton-Cotes.}]
class NewtonCotes:
    @staticmethod
    def _composite(func, a, b, n, rule):
        h = (b - a) / n
        return sum(rule(func, a + i * h, a + (i + 1) * h) for i in range(n))
\end{lstlisting}

Por exemplo, aplicando-se a regra de Simpson $1/3$ à função de teste $f(x) = (\sin(2x) + 4x^2 + 3x)^2$ em $[0, 1]$, cuja resposta analítica rigorosa converge para $\approx 17.8764703$, obtemos os seguintes resíduos práticos em execução:
\begin{itemize}
    \item \textbf{Trapézio} (100 partições): $17.877811$
    \item \textbf{Simpson 1/3} (10 partições): $17.876473$
    \item \textbf{Regra de Boole} (5 partições): $17.876470$
\end{itemize}
Fica cristalino que a elevação do grau subjacente ao interpolador (na regra de Boole, utiliza-se $n=4$) drásticamente abate a necessidade de excessivas partições, proporcionando estabilidade imediata.

\section{Quadratura Gaussiana}
As limitações atreladas à imposição de malhas unifomemente espaçadas deram lugar às otimizações de Carl Friedrich Gauss. Ao permitirmos que as posições dos nós ($x_i$) e os respectivos pesos numéricos ($w_i$) entrem como variáveis estritamente livres num sistema polinomial, uma quadratura de grau $N$ é capaz de integrar sem nenhum erro intrínseco polinômios de até grau $2N-1$.

Na biblioteca desenvolvida, a resolução das raízes e os pesos dinâmicos foram instanciados recorrendo às funções acopladas em \texttt{numpy.polynomial}, proporcionando suporte paramétrico a múltiplos métodos:
\begin{itemize}
    \item \textbf{Gauss-Legendre}: Utilizado para $I \in [a, b]$ clássicos, transformando-os biunivocamente em $[-1, 1]$.
    \item \textbf{Gauss-Hermite}: Adequado a limites infinitos e subjacentes a funções estatísticas do tipo $f(x)e^{-x^2}$.
    \item \textbf{Gauss-Laguerre}: Voltado a distribuições de decaimento exponencial $f(x)e^{-x}$ de $[0, \infty]$.
    \item \textbf{Gauss-Chebyshev}: Com pesos imutáveis $w = \frac{\pi}{n}$, domina integrais afetadas por denominador em forma de arco.
\end{itemize}

O desempenho destes estimadores é notório. A integração de Gauss-Legendre de grau $n=4$ para o mesmo polinômio de teste supracitado retornou exatamente $17.876480$, alcançando o platô com avaliações consideravelmente inferiores às requeridas em Newton-Cotes Composta.

\section{Tratamento Analítico de Singularidades}
Se a função divergir vertiginosamente nos extremos da integral ($f(x) \to \infty$ quando $x \to 0$, por exemplo), os métodos fechados clássicos geram erro aritmético catastrófico de \texttt{DivisionByZero}. 
O escape teórico é o mapeamento hiperbólico. Através do Jacobiano de transformações não-lineares, "esticamos" o espaço do infinito concentrando amostras densamente em volta dos polos. 

Implementamos o método da \textbf{Exponencial Simples} via tangente hiperbólica ($x(t) \propto \tanh(t)$) e a respectiva evolução, a \textbf{Duplo Exponencial}, baseada na superposição de senos hiperbólicos ($x(t) \propto \tanh(\sinh(t))$).

Ao aplicarmos as rotinas à integral singular $\int_0^1 x^{-1/2}dx = 2$, os algoritmos provaram sua eficácia evadindo a origem pura e computando o limite lateral com erro residual de alguns pontos percentuais, mesmo cortando deliberadamente as caudas em um $T=3.5$.

\begin{lstlisting}[language=Python, caption={Trecho: Transformação Exponencial Simples.}]
class ExponentialTransforms:
    @staticmethod
    def simple_exponential(f, a, b, n, T=3.5):
        h = 2 * T / n
        t = np.linspace(-T, T, n+1)
        phi = np.tanh(t)
        dphi = 1.0 / np.cosh(t)**2
        # Mapeia para [a,b]
        x = 0.5 * (a + b) + 0.5 * (b - a) * phi
        # ...
\end{lstlisting}

\section{Integração Multidimensional}
A transição das teorias fundamentais para volumes delimitados requer somas iteradas de Riemann. A arquitetura adotada parametrizou a varredura ortogonal do espaço integrando inicialmente em relação a $x$ com pesos isolados e iterando o resultado na superposição perante a dimensão $y$. 

A função 2D acoplada foi implementada com uma Regra de Simpson Bidimensional otimizada. Sob $f(x,y) = xy$ integrando em matriz $1 \times 1$, o retorno $0.250000$ bate precisamente com o modelo teórico de $1/4$.

\section{Conclusão}
A exploração sistêmica promovida na Unidade 2 certifica a absoluta relevância de transitar por diferentes bases de cálculo numérico dependendo das injunções do problema. Fórmulas de Newton-Cotes apresentam grande resiliência e facilidade interpretativa para curvas comportadas; entretanto, integrandos munidos de alta oscilação ou crescimento estocástico em limites finitos clamam veementemente pela densidade das raízes ortogonais de Gauss, ou sob a iminência de polos assintóticos, pelas transformações duplas hiperbólicas.

As respostas obtidas referendam de modo irrefutável a estabilidade dos módulos \textit{Python} desenvolvidos. O avanço metodológico oriundo da vetorização e das rotinas nativas abrigou as antigas dezenas de laços interativos em curtas invocações diretas de matrizes NumPy.

\end{document}
